\documentclass[12pt,a4paper]{article}
% ============================================================
%  Structural Persistence Series — Shared Preamble
% ============================================================
\usepackage{fontspec}
\setmainfont{Hiragino Mincho ProN}
\setsansfont{Hiragino Sans}
\setmonofont{Menlo}

\usepackage{iftex}
\ifXeTeX
  \XeTeXlinebreaklocale "ja"
  \XeTeXlinebreakskip=0pt plus 1pt minus 0.1pt
\fi

\usepackage[margin=22mm]{geometry}
\usepackage{amsmath,amssymb,amsthm}
\usepackage{xcolor}
\usepackage{graphicx}
\usepackage{hyperref}
\usepackage{booktabs}
\usepackage{fancyhdr}
% enumitem not available in this TeX installation

% --- Colours ------------------------------------------------
\definecolor{PaperAccent}{HTML}{1F4E79}
\definecolor{PaperGray}{HTML}{51606F}

\hypersetup{
  colorlinks=true,
  linkcolor=PaperAccent,
  urlcolor=PaperAccent,
  citecolor=PaperAccent,
  pdflang=ja
}
\urlstyle{same}

% --- Typography ---------------------------------------------
\setlength{\parindent}{1.5em}
\setlength{\parskip}{0pt}
\setlength{\emergencystretch}{3em}
\linespread{1.08}
\setlength{\abovedisplayskip}{8pt plus 2pt minus 4pt}
\setlength{\belowdisplayskip}{8pt plus 2pt minus 4pt}
\setlength{\abovedisplayshortskip}{6pt plus 2pt minus 2pt}
\setlength{\belowdisplayshortskip}{6pt plus 2pt minus 2pt}

% --- Section label names ------------------------------------
\renewcommand{\abstractname}{要旨}

% --- Theorem-like environments ------------------------------
\theoremstyle{plain}
\newtheorem{proposition}{命題}[section]
\newtheorem{lemma}[proposition]{補題}
\newtheorem{corollary}[proposition]{系}

\theoremstyle{definition}
\newtheorem{definition}{定義}[section]
\newtheorem{assumption}{条件}[section]

\theoremstyle{remark}
\newtheorem*{remark}{注}

% --- Running header -----------------------------------------
\pagestyle{fancy}
\fancyhf{}
\renewcommand{\headrulewidth}{0.4pt}
\fancyhead[L]{\small\textcolor{PaperGray}{\nouppercase{\leftmark}}}
\fancyhead[R]{\small\textcolor{PaperGray}{\thepage}}
\fancypagestyle{plain}{%
  \fancyhf{}
  \renewcommand{\headrulewidth}{0pt}
  \fancyfoot[C]{\small\textcolor{PaperGray}{\thepage}}
}

% --- Title block macro --------------------------------------
\newcommand{\PaperTitleBlock}[3]{%
  % #1 = title, #2 = subtitle, #3 = date
  \thispagestyle{plain}%
  \begin{center}
  {\LARGE\bfseries #1\par}
  \vspace{0.4em}
  {\normalsize #2\par}
  \vspace{1.0em}
  {\large Akihito Sunagawa\par}
  \vspace{0.3em}
  {\normalsize #3\par}
  \end{center}
  \vspace{1.6em}
}

% Legacy 2-arg version (backward compat)
\newcommand{\PaperTitleBlockSimple}[2]{%
  \PaperTitleBlock{#1}{}{#2}
}


\begin{document}

\PaperTitleBlock
  {構造持続の最小核と収支原理}
  {— 維持可能領域の縮小、対数指数核、回復を含む持続条件 —}
  {2026年4月12日}

\begin{abstract}
本稿は、分冊版である「構造持続の最小形式」と「構造持続の収支原理」を、外部読者向けに一つの導線として再構成した統合短論文である。新しい理論核、新しい定理、新しい実験的 support を追加するものではない。厳密な公理、定理、証明、限界条件は、分冊版である Paper 1 / Paper 2 および対応する技術補論を基準とする。

ここでいう「最小」は、当たり前という意味ではない。Paper 1 の最小形式は、構造の喪失を資源枯渇ではなく維持可能領域の縮小として捉え、対数比尺度を一意化し、事後的表現選択による空虚化を防ぐ条件を切り出す、非自明な表現定理である。Paper 2 は、その最小形式を前提として、回復を含む系へ拡張する。

構造が失われるとは、存在一般が消えることではなく、その構造を維持できる状態の領域が縮小することである。事前固定された構造維持問題において、維持可能領域の残存比を連鎖的に測る尺度は、自然な合成性条件のもとで対数比として一意に定まる。その結果、最小核は

\[
S = M e^{-L}
\]

の形を取る。ここで \(L\) は累積構造消耗量、\(M\) は有効維持資源である。

回復、修復、再構成、外部支援を明示する場合には、各時点の構造消耗量を \(d_{t}\)、回復量を \(r_{t}\) とし、純消耗量

\[
b_{t}=d_{t}-r_{t}
\]

を導入する。累積純消耗量

\[
B_{n}=\sum_{t<n} b_{t}
\]

により、回復を含む構造持続量は

\[
S_{n}=M_{n} e^{-B_{n}}
\]

と読める。本稿の中心は、閉じた縮小の最小核 \(S=Me^{-L}\) が、開いた系の収支原理 \(S=Me^{-B}\) へ自然に拡張されることを、読者向けに一つの流れとして示すことである。
\end{abstract}

\section{本稿の位置づけ}\label{ux672cux7a3fux306eux4f4dux7f6eux3065ux3051}

本稿は、主理論 spine の統合読解版である。役割は、Paper 1 と Paper 2 の内容を一つの短い道筋として読ませることにある。

{\def\LTcaptype{table} % do not increment counter
\begin{longtable}[]{@{}
  >{\raggedright\arraybackslash}p{(\linewidth - 2\tabcolsep) * \real{0.5000}}
  >{\raggedright\arraybackslash}p{(\linewidth - 2\tabcolsep) * \real{0.5000}}@{}}
\toprule\noalign{}
\begin{minipage}[b]{\linewidth}\raggedright
文書
\end{minipage} & \begin{minipage}[b]{\linewidth}\raggedright
役割
\end{minipage} \\
\midrule\noalign{}
\endhead
\bottomrule\noalign{}
\endlastfoot
Paper 0 & 理論体系全体の地図、観測可能性の三層、実証・Lean・補論群の位置づけ \\
本 Core Paper & Paper 1 / Paper 2 を外部読者向けに一本の導線として読む統合短論文 \\
Paper 1 & 構造消耗尺度、対数比の一意性、望遠鏡積、最小核の厳密 spine \\
Paper 2 & 回復を含む二段階更新、純消耗量、収支原理の厳密 spine \\
\end{longtable}
}

したがって、本稿で述べる主張の厳密な基準は分冊版にある。本稿は読者導線であり、分冊版の置き換えではない。

とくに、Paper 1 は単なる予備的な自明事項ではない。構造を「その構造として維持できる状態集合」として固定し、その縮小を測る尺度がなぜ対数比でなければならないかを示すことが、主理論の第一の非自明な仕事である。

\section{問い}\label{ux554fux3044}

長期運用系、学習系、推論系、制約充足系では、資源が完全に尽きていなくても構造が保てなくなることがある。このとき問うべきなのは、単に「どれだけ資源が残っているか」ではない。

問うべきなのは、ある構造を維持できる状態の領域が、制約、矛盾、損傷、更新、外乱、修復のもとでどのように変化するかである。

この問いを最小化すると、次の二段階になる。

\begin{enumerate}
\def\labelenumi{\arabic{enumi}.}
\tightlist
\item
  回復を明示しない閉じた縮小では、構造消耗はどの形で測られるべきか。
\item
  回復を含む開いた系では、構造持続は何によって支配されるか。
\end{enumerate}

Paper 1 は第一の問いに答え、Paper 2 は第二の問いに答える。本稿は、その二つを一続きの物語として読む。

\section{事前固定された構造維持問題}\label{ux4e8bux524dux56faux5b9aux3055ux308cux305fux69cbux9020ux7dadux6301ux554fux984c}

系 \(X\) と、構造を維持できる状態集合 \(V\) を考える。初期の維持可能領域を \(V^{(0)}\) とし、制約の蓄積により

\[
V^{(0)} \supseteq V^{(1)} \supseteq \cdots \supseteq V^{(n)}
\]

と縮小するとする。

ここで比較に用いる測度 \(m\)、対象構造、段階列、時間地平は事前に固定されていなければならない。観測後に \(V\) や \(m\) を選び直せば、形式は空虚化する。これは Paper 1 の中心的な防御線である。

以下で扱う比は、比較する集合の質量が正かつ有限である範囲で読む。ゼロ到達は通常の対数比ではなく、collapse boundary として別扱いする。

\section{構造消耗尺度と対数比}\label{ux69cbux9020ux6d88ux8017ux5c3aux5ea6ux3068ux5bfeux6570ux6bd4}

残存比率 \(r\in(0,1]\) に対する構造消耗尺度 \(f(r)\) を考える。尺度に求めるのは、制約生成過程の独立性ではない。求めるのは、すでに得られた残存比を連鎖的に測るときの合成性である。

すなわち、

\[
\begin{aligned}
\frac{m(V^{(2)})}{m(V^{(0)})} \\
&= \\
\frac{m(V^{(1)})}{m(V^{(0)})} \\
\frac{m(V^{(2)})}{m(V^{(1)})}
\end{aligned}
\]

と分解されるなら、対応する消耗量は和として合成されるべきである。

この要請は、制約が独立に効くという経験的仮定ではない。制約どうしの重なり、依存、冗長性は、各段階の \(V^{(i)}\) と測度比

\[
\frac{m(V^{(i)})}{m(V^{(i-1)})}
\]

にすでに反映される。

この合成性、正規化、連続性のもとで、構造消耗尺度は

\[
f(r)=-k\log r
\]

と表される。単位規約として \(k=1\) を取れば、段階消耗量は

\[
\begin{aligned}
d_{i} \\
&= \\
-\log\frac{m(V^{(i)})}{m(V^{(i-1)})}
\end{aligned}
\]

である。

\section{望遠鏡積と最小核}\label{ux671bux9060ux93e1ux7a4dux3068ux6700ux5c0fux6838}

累積構造消耗量を

\[
L_{n}=\sum_{i=1}^{n} d_{i}
\]

と置くと、望遠鏡積により

\[
m(V^{(n)})=m(V^{(0)})e^{-L_{n}}
\]

が恒等的に成り立つ。

これは「世界が経験的に指数減衰する」という主張ではない。事前固定された \(V,m\) の残存比を対数尺度で測ると、望遠鏡積から指数核が現れる、という表現定理である。望遠鏡積は座標を採用した後の代数的帰結であり、本稿の非自明な部分は、構造の失敗を維持可能領域の縮小として読む座標そのものを切り出す点にある。

ここに有効維持資源 \(M\) を別途導入すると、最小の構造持続量は

\[
S=Me^{-L}
\]

と書ける。\(M\) は指数核そのものから導かれる量ではなく、構造を維持するために実際に使える資源、余力、支援能力を表す外側の資源項である。

\section{回復を含む二段階更新}\label{ux56deux5fa9ux3092ux542bux3080ux4e8cux6bb5ux968eux66f4ux65b0}

開いた系では、構造維持可能領域は単に縮小するだけではない。修復、回復、冗長化、外部支援、ロールバック、再構成によって、維持可能領域が再拡大する場合がある。

そこで一つの観測単位の中で、まず収縮後の中間集合 \(V_{t}^-\) を置き、その後に回復後の集合 \(V^{(t+1)}\) を置く。

構造消耗量を

\[
\begin{aligned}
d_{t} \\
&= \\
-\log\frac{m(V_{t}^-)}{m(V^{(t)})}
\end{aligned}
\]

回復量を

\[
\begin{aligned}
r_{t} \\
&= \\
\log\frac{m(V^{(t+1)})}{m(V_{t}^-)}
\end{aligned}
\]

と定義すると、純消耗量は

\[
\begin{aligned}
b_{t}&=d_{t}-r_{t} \\
&= \\
-\log\frac{m(V^{(t+1)})}{m(V^{(t)})}
\end{aligned}
\]

となる。中間集合 \(V_{t}^-\) は差し引きの中で打ち消し合う。

ただし、これは endpoint が与えられた後の対数比計算に関する事実である。実際の力学として収縮作用と回復作用の順序を変えれば、終点 \(V^{(t+1)}\) 自体が変わる可能性がある。したがって、形式的な打ち消しを、実体的な順序不変性と読んではならない。

\section{収支原理}\label{ux53ceux652fux539fux7406}

有限時間地平 \(n\) までの累積純消耗量を

\[
B_{n}=\sum_{t<n} b_{t}
\]

と置くと、

\[
m(V^{(n)})=m(V^{(0)})e^{-B_{n}}
\]

が成り立つ。

資源項を含めれば、回復を含む構造持続量は

\[
S_{n}=M_{n} e^{-B_{n}}
\]

である。

これは最小核を置き換えるものではない。回復がない場合、\(r_{t}=0\) であり、\(B_{n}=L_{n}\) に戻る。したがって、Paper 2 は Paper 1 を否定するのではなく、Paper 1 の最小核を、回復を含む開いた系へ拡張する。

\section{収支は均衡ではない}\label{ux53ceux652fux306fux5747ux8861ux3067ux306fux306aux3044}

本稿でいう収支は equilibrium を意味しない。消耗、維持、回復を同じ差し引き量の符号で読むための accounting を意味する。

{\def\LTcaptype{table} % do not increment counter
\begin{longtable}[]{@{}ll@{}}
\toprule\noalign{}
条件 & 読み方 \\
\midrule\noalign{}
\endhead
\bottomrule\noalign{}
\endlastfoot
(b\_t\textgreater0) & 消耗優位 \\
(b\_t=0) & 維持局面 \\
(b\_t\textless0) & 回復優位 \\
\end{longtable}
}

この三局面は、系が均衡状態にあるかどうかの判定ではない。事前固定された観測単位の上で、差し引き量の符号を読むための局所的な会計分類である。

隣接する二つの区間を併合すれば

\[
\begin{aligned}
b_{t:t+2} \\
&= \\
-\log\frac{m(V^{(t+2)})}{m(V^{(t)})} \\
&= \\
b_{t}+b_{t+1}
\end{aligned}
\]

となるため、累積純消耗量 \(B_{n}\) は時間併合に対して加法的に整合する。一方で、細かい粒度では「消耗の後に回復」と読める経路が、粗い粒度では \(b_{t:t+2}=0\) の維持局面として読まれることがある。

したがって、局面分類、collapse boundary、hitting-time bound を経験的または確率的に主張する場合には、対象構造だけでなく、観測単位、段階列、時間地平も事前に固定する必要がある。

\section{観測可能性の三層}\label{ux89b3ux6e2cux53efux80fdux6027ux306eux4e09ux5c64}

同じ構造持続核は、観測可能性の異なる三つの層で読まれる。

{\def\LTcaptype{table} % do not increment counter
\begin{longtable}[]{@{}
  >{\raggedright\arraybackslash}p{(\linewidth - 2\tabcolsep) * \real{0.5000}}
  >{\raggedright\arraybackslash}p{(\linewidth - 2\tabcolsep) * \real{0.5000}}@{}}
\toprule\noalign{}
\begin{minipage}[b]{\linewidth}\raggedright
層
\end{minipage} & \begin{minipage}[b]{\linewidth}\raggedright
読み方
\end{minipage} \\
\midrule\noalign{}
\endhead
\bottomrule\noalign{}
\endlastfoot
仕様固定構造層 & 構造、測度、境界を仕様から直接固定できる \\
条件付き構造埋め込み層 & 既存理論の drift、差分、停止境界を本理論の変数へ条件付きに写す \\
構造推定層 & 構造そのものを直接数えず、代理指標と凍結検証で推定する \\
\end{longtable}
}

この三層は、理論の強弱や対象ドメインの固定分類ではない。違うのは、\(V,m,d_{t},r_{t},M\) をどの程度直接指定できるかである。

仕様固定構造層では、定理、有限時間境界、限定クラス統一 interface が主役になる。構造推定層では、proxy、frozen validation、baseline + structural persistence の追加予測力が主役になる。条件付き構造埋め込み層では、既存理論との formal bridge が主役になる。

\section{外向け価値}\label{ux5916ux5411ux3051ux4fa1ux5024}

この統合読解版で見せたい価値は、単に二つの式を並べることではない。第一に、\(L\) と \(B\) は、自然対数で測られる無次元の対数比であり、事前固定された写像のもとで構造消耗と回復を定量的に読むための共通座標になる。単位規約 \(k=1\) を取ると、この座標は自然対数単位、すなわち nats で読める。

第二に、経験的価値は、構造持続座標だけが常に最強の単独モデルになることではない。多くの現実ドメインには、既存専門モデルや強い domain baseline がすでにある。したがって中心的な検査は、構造持続座標を追加した

\[
\text{domain baseline}+\text{SP}
\]

が、凍結された out-of-sample 条件で domain baseline 単独を改善するかである。ここで SP は structural persistence coordinate、すなわち構造消耗、回復、余力、代替経路などの指標群を指す。

第三に、この座標は設計転用の言語になる。あるドメインで効いた consumption localization、recovery preservation、margin preservation、alternative-path preservation は、別ドメインの candidate intervention を作る手がかりになる。ただし、転用されるのは support ではない。別ドメインでの support は、そのドメインで写像、指標、baseline、metric、split、判定規則を凍結し、holdout、future surface、fresh archive、または outside rerun で検証して初めて成立する。

\section{本稿で言っていないこと}\label{ux672cux7a3fux3067ux8a00ux3063ux3066ux3044ux306aux3044ux3053ux3068}

本稿は、すべての系が指数減衰するという経験法則を主張しない。

本稿は、任意のドメインに対する単一普遍不等式を主張しない。

本稿は、構造推定層の proxy が仕様固定構造層と同じ強さの証拠を持つとは主張しない。

本稿は、\(M\) が指数核から導かれるとは主張しない。\(M\) は有効維持資源を表す外側の資源項であり、その操作化は別補論の対象である。

本稿が主張するのは、事前固定された構造維持問題において、構造消耗を対数比として測ると最小核 \(S=Me^{-L}\) が現れ、回復を同じ尺度で差し引くと収支核 \(S=Me^{-B}\) が得られる、という一点である。

\section{結論}\label{ux7d50ux8ad6}

構造持続理論の最小核は、維持可能領域の縮小を対数比で測ることから現れる。

回復を含む開いた系では、構造持続を決めるのは消耗量そのものではなく、消耗量と回復量の差である。

したがって、外部読者向けの最短導線は次の一本線である。

\[
\begin{aligned}
S&=Me^{-L} \\
\quad\longrightarrow\quad \\
S&=Me^{-B}, \\
\qquad \\
B_{n}&=\sum_{t<n}(d_{t}-r_{t}).
\end{aligned}
\]

この一本線を厳密に支える分冊版が Paper 1 と Paper 2 であり、本稿はその統合読解版である。
\end{document}
